% Generated draft: requires mathematical and source review.
\documentclass[11pt]{article}
\usepackage[margin=25mm]{geometry}
\usepackage{fontspec}
\setmainfont{Latin Modern Roman}
\usepackage{amsmath,amssymb,mathtools,unicode-math}
\setmathfont{Latin Modern Math}
\usepackage{enumitem,needspace,xurl,hyperref}
\setlength{\parindent}{0pt}
\setlength{\parskip}{5pt}
\setlength{\emergencystretch}{3em}
\newcommand{\N}{\mathbb{N}}
\newcommand{\Z}{\mathbb{Z}}
\newcommand{\Q}{\mathbb{Q}}
\newcommand{\R}{\mathbb{R}}
\newcommand{\C}{\mathbb{C}}
\newcommand{\F}{\mathbb{F}}
\newcommand{\A}{\mathbb{A}}
\newcommand{\PP}{\mathbb{P}}
\newcommand{\E}{\mathbb{E}}
\newcommand{\eps}{\varepsilon}
\newcommand{\dd}{\,\mathrm{d}}
\begin{document}
\section*{UnsolvedMath ID: 11000058}
\textbf{Source problem code:} catalogue entry
\subsection{Definitions and mathematical statement}
In the standard closed-surface convention, let $S_g$ be a closed connected oriented surface of genus $g$ and let
\[
 \mathcal I_g=\ker\bigl(\operatorname{Mod}(S_g)
 \to\operatorname{Sp}(H_1(S_g;\Z))\bigr)
\]
be its Torelli group. Conjecture 5.12 states
\[
 \forall g\geq4,\qquad\mathcal I_g\text{ is finitely presented}.
\]
That is, for every such $g$, there should be a presentation $\mathcal I_g\cong\langle x_1,\ldots,x_m\mid r_1,\ldots,r_s\rangle$ with both $m,s$ finite. Equivalently, $\mathcal I_g$ should have a $K(\mathcal I_g,1)$ with finite $2$-skeleton. The conjecture's explicit threshold is $g\geq4$, not $g\geq3$.

The source contrasts this with a statement that neither $\mathcal I_g$ nor $K_g$ has a classifying space homotopy equivalent to a finite complex, and announces a stronger result due to Akita. The excerpt does not include that result, its genus restrictions, or the definition of $K_g$ (customarily a Johnson kernel). These background claims therefore require source confirmation in their precise range. Absence of a finite classifying complex does not itself preclude a finite $2$-skeleton.

The registry reports finite generation for $g\geq3$ and attributes the finite-presentability question to Mess and Kirby's problem list, while reporting no resolution through 2026. This is supplied status information, not an independent modern assessment.
\subsection{Short English statement}
For every genus $g\geq4$, the Torelli group should admit a presentation with finitely many generators and relations. This asks for a finite $2$-skeleton, not a finite entire classifying space.
\subsection{Sources}
\begin{itemize}
\item Local registry: \url{data/unsolvedmath/problems.json}, permanent UnsolvedMath ID $11000058$, exact source code \texttt{catalogue entry}.
\item Farb, \emph{Problems on Mapping Class Groups and Related Topics}, Conjecture 5.12, PDF page 47: \url{https://www.math.uchicago.edu/~farb/papers/mcgbook.pdf}. Supplied access date: 2026-07-29; extraction: pdf-font-structure; page-disambiguated numbering among 329 headings.
\item Supplied attribution: Mess's conjecture/problem, Kirby's list Problem 2.5(A); cross-reference to another chapter's Problem 3.1, source code \texttt{AMR-109-0305}. Full references are absent.
\item The excerpt attributes a following stronger nonfiniteness result to Akita, but the theorem itself and its bibliography are not supplied.
\end{itemize}
\end{document}
